\documentclass[english,article,nodocrefonfront,nodocrefonback]{utbmciadreport}

\usepackage[T1]{fontenc}

\declaredocument{UTBM/CIAD Report Template Extension to TeX-UPmethodology}{Official Documentation - Version \insertciadreportthemeversion}{utbmciad-report-doc}

\addauthorvalidator[galland@arakhne.org]{St{\'e}phane}{Galland}

\setdockeywords{utbmciad, CIAD, UTBM, Style Extension, TeX-UPmethodology, \LaTeX}

\setdocabstract{This document describes the template extension to TeX-UPmethodology which is providing the standard document layout and colors of the Laboratory ``Connaissance et Intelligence Artificielle Distribuées'' from the Université de Technologie Belfort-Montbéliard, France.}

\updateversion{1.0}{\makedate{05}{08}{2019}}{First version of the documentation.}{\upmpublic}
\updateversion{1.1}{\makedate{09}{01}{2025}}{Remove UBFC.}{\upmpublic}
\updateversion{1.2}{\makedate{06}{04}{2026}}{Add funding scheme features.}{\upmpublic}
\updateversion{1.3}{\makedate{10}{04}{2026}}{Add color for watermark.}{\upmpublic}
\incsubversion{\makedate{23}{04}{2026}}{Add UTF8 support package.}{\upmpublic}

%-------------------------------
% CAUTION: State the capture of the definitions of the theme. YOU MUST NOT USE THESE COMMANDS IN YOUR DOCUMENT.
%-------------------------------

\makeatletter
\upm@extension@savealllangtrue

% The following commands are defined to capture the definitions of the colors and the values.
\gdef\upmextensiondoccolorlist{}
\gdef\upmextensiondocvaluelist{}

\global\let\upm@extensiondoc@old@definecolor\definecolor
\global\let\upm@extensiondoc@old@setpdfcolor\setpdfcolor
\global\let\upm@extensiondoc@old@setvalue\upm@extension@Set

\renewcommand{\definecolor}[3]{%
	\upm@extensiondoc@old@definecolor{#1}{#2}{#3}%
	\protected@xdef\upmextensiondoccolorlist{\upmextensiondoccolorlist%
		\protect\ensuremath{\protect\upm@protect{#1}} & #2 & #3 & \protect\textcolor{#1}{\protect\ensuremath{\protect\upm@protect{#1}}} \protect\\%
	}%
}

\renewcommand{\setpdfcolor}[1]{%
	\upm@extensiondoc@old@setpdfcolor{#1}%
	\upm@extensiondoc@old@definecolor{my__pdf__color}{rgb}{#1}%
	\protected@xdef\upmextensiondoccolorlist{\upmextensiondoccolorlist%
		``PDFCOLOR'' & rgb & #1 & \protect\textcolor{my__pdf__color}{PDF color} \protect\\%
	}%
}

\let\MyProtect\upm@protect
\newcommand{\MyGet}[2]{%
	\ifthenelse{\equal{#2}{copyright}}{-}{%
	\ifthenelse{\equal{#2}{trademarks}}{-}{%
	\ifthenelse{\equal{#2}{cfrontpage}}{-}{%
	\ifthenelse{\equal{#2}{backpage}}{-}{%
	\ifthenelse{\equal{#2}{abstract}}{-}{%
	\ifthenelse{\equal{#2}{keywords}}{-}{%
	\ifthenelse{\equal{#2}{logo}}{-}{%
	\ifthenelse{\equal{#2}{smalllogo}}{-}{%
	\ifthenelse{\equal{#2}{frontillustration}}{-}{%
	\ifthenelse{\equal{#2}{headerillustration}}{-}{%
	\GetLang{#2}{#1}%
	}}}}}}}}}}%
}

\renewcommand{\Set}[3][default]{%
	\ifthenelse{\equal{#1}{default}}{%
		\upm@extensiondoc@old@setvalue{#2}{#3}%
		\protected@xdef\upmextensiondocvaluelist{\upmextensiondocvaluelist%
			\protect\upm@protect{#2} & \protect\textit{all} & \protect\MyGet{\upmcurrentlang}{#2} \protect\\%
		}%
	}{%
		\upm@extensiondoc@old@setvalue[#1]{#2}{#3}%
		\protected@xdef\upmextensiondocvaluelist{\upmextensiondocvaluelist%
			\protect\upm@protect{#2} & \protect\upm@protect{#1} & \protect\MyGet{#1}{#2} \protect\\%
		}%
	}%
}

\input{upmext-utbmciad.cfg}

\makeatother
%-------------------------------
% END OF SECTION
%-------------------------------

\begin{document}

\section{Introduction}

This document describes the style extension for TeX-UPmethodology\footnote{\url{https://www.arakhne.org/latex/tex-upmethodology}} which is providing the standard document layout and colors of the Laboratory ``Connaissance et Intelligence Artificielle Distribuées''\footnote{\url{https://www.ciad-lab.fr}} (CIAD UR 7533) of ``Université de Technology Belfort-Monbéliard''\footnote{\url{https://www.utbm.fr}}, France.

\Emph{This document contains only the values and the macros provided by the extensions itself, not the values and macros provided by TeX-UPmethodology.}

\section{Usage and Licensing}

This extension contains two distinct types of material:
\begin{enumerate}
\item LaTeX source code (all files that are under the \texttt{tex/} and \texttt{docs/} directories and any documentation source files).
\item University and institutions logo images (all files under the \texttt{logos/} directory).
\end{enumerate}

Licensing terms for each part are different:
\begin{enumdescription}
\item[LaTeX source code] The LaTeX source code is free software: you can redistribute it and/or modify it under the terms of the GNU Lesser General Public License as published by the Free Software Foundation, either version 3 of the License, or (at your option) any later version.

	A copy of the GNU Lesser General Public License (LGPLv3) is provided in the file \texttt{LICENSE.LGPLv3}.  You can also find it at \url{https://www.gnu.org/licenses/lgpl-3.0.html}.
\item[University and institution logos] The logo files are not licensed under the LGPLv3. They are the property of their respective university or institution and are included with permission for redistribution only as part of this complete package. They may NOT be modified, altered, or extracted and used separately without explicit written permission from the owning university or institution.

	The exact terms governing the logos are set forth in the file \texttt{logos/LICENSE.logos}. By redistributing this package, you agree to abide by those terms.
\end{enumdescription}

When redistributing this package (e.g., on CTAN, a personal website, or a CDN), you must preserve the directory structure and all accompanying license files.  Users of the LaTeX code may remove the \texttt{logos/} directory without affecting their rights under the LGPLv3.

\section{Requirements}

\texttt{utbmciadreport} requires the core packages of \texttt{tex-upmethodology}.

\section{Installation}

Copy \texttt{utbmciad-report} configuration file (\texttt{.cfg}) and all related pictures in your system-wide or local user \texttt{texmf}.

\section{Class Options}

\texttt{utbmciad-report} provides several class options that are specific to the theme.
The other class options are passed to \texttt{tex-upmethodology}.


\subsection{Language}

\begin{description}
\item[english] The document is written in english (default).
\item[french] The document is written in french.
\end{description}

\subsection{Logos}

\begin{description}
\item[ciadlogo] Show up the CIAD logo on the front and back pages.
\item[utbmlogo] Show up the UTBM logo on the front and back pages.
\item[noutbmlogo] Hide the UTBM logo from the front and back pages.
\item[utlogo] Show up the UT Network logo on the front and back pages.
\item[noutlogo] Hide the UT Network logo from the front and back pages.
\item[nociadlogo] Hide the CIAD logo from the front and back pages.
\item[icartslogo] Show up the iC ARTS logo on the front and back pages.
\item[noicartslogo] Hide the iC ARTS logo from the front and back pages.
\item[sayenslogo] Show up the SATT Sayens logo on the front and back pages.
\item[nosayenslogo] Hide the SATT Sayens logo from the front and back pages.
\item[fundingschemelogo] Show up the logos of the funding schemes.
\item[nofundingschemelogo] Hide the logos of the funding schemes.
\item[doclogo] Same as \texttt{ciadlogo} option.
\item[nodoclogo] Same as \texttt{nociadlogo} option.
\end{description}

\subsection{Document Reference}

\begin{description}
\item[docrefonfront] Show up the document reference on the front page.
\item[nodocrefonfront] Hide the document reference from the front page.
\item[docrefonback] Show up the document reference on the back page.
\item[nodocrefonback] Hide the document reference from the back page.
\end{description}

\subsection{Contact Information}

\begin{description}
\item[contactonback] Show up the contact information on the back page.
\item[nocontactonback] Hide the contact information from the back page.
\end{description}

\section{Provided Macros}

\texttt{utbmciad-report} provides the following macros:
\begin{itemize}
\item \texttt{{\textbackslash}addfundingschemelogo[height size]\{path to logo image\}}: Add a logo image in the list of the funding schemes' logos. The default height is $1.25cm$.
\item \texttt{{\textbackslash}utbmciadfundingschemeslogos}: Show the list of all the funding schemes' logos.
\end{itemize}

The packages of \texttt{tex-upmethodology} (a dependency of this style) provides plenty of additional macros. See the documentation of \texttt{tex-upmethodology} for details.

\section{Provided Colors}

\texttt{utbmciad-report} provides the following colors. The colors are expressed according to the standard \texttt{color.sty} specifications. If the table below is empty, it means that this extension does not provide any specific color.

\begin{mtabular}{4}{XccX}{Provided Colors}
\tabularheader{Name}{Type}{Value}{Example}
\upmextensiondoccolorlist
\end{mtabular}

\section{Provided Values}

\texttt{utbmciad-report} provides the following values. The values are accessible with \texttt{{\textbackslash}Get} macro. If the table below is empty, it means that this extension does not provide any specific definition of value.

\begin{tiny}
\begin{mtabular}{3}{XcX}{Provided Values}
\tabularheader{Name}{Language}{Content}
\upmextensiondocvaluelist
\end{mtabular}
\end{tiny}

\end{document}
